\section{Informations complémentaires}

\subsection{Tableau complet des connexions}
\begin{table}[H]
\centering
\begin{tabular}{|l|l|l|l|}
\hline
\textbf{Composant} & \textbf{Pin} & \textbf{Fonction} & \textbf{Configuration} \\
\hline
HC-SR04 Trigger & PB7 & Output & TIM4\_CH2 (PWM) \\
\hline
HC-SR04 Echo & PB8 & Input & EXTI (Rising/Falling) \\
\hline
Servo & PB5 & Output & TIM3\_CH2 (PWM) \\
\hline
LED Bleue & PD15 & Output & GPIO Push-Pull \\
\hline
LED Verte & PD12 & Output & GPIO Push-Pull \\
\hline
UART TX & PA2 & Output & USART2\_TX \\
\hline
UART RX & PA3 & Input & USART2\_RX \\
\hline
Bouton-poussoir & PA0 & Input & EXTI (Falling) \\
\hline
\end{tabular}
\caption{Récapitulatif des connexions matérielles}
\label{tab:connections}
\end{table}

\subsection{Gestion des erreurs}
Le système intègre une gestion basique des erreurs:
\begin{itemize}
    \item Vérification des mesures de distance (filtrage des valeurs aberrantes)
    \item Détection des timeouts pour le capteur ultrasonique
    \item Validation des commandes reçues via la liaison série
\end{itemize}

\subsection{Améliorations possibles}
Plusieurs améliorations pourraient être apportées:
\begin{itemize}
    \item Implémentation d'un filtre pour stabiliser les mesures de distance
    \item Développement d'une interface utilisateur plus complète
    \item Ajout de fonctionnalités de diagnostic et de calibration
    \item Utilisation du DMA pour optimiser les transferts UART
\end{itemize}